% ===========================================================================
\chapter{Working out everything}\label{ch:fixpoint}

One application of a rule is not enough. Applying the rules for
\code{sees} once finds whose posts reach Ann directly, and through one
reshare, but a post three reshares away needs the two-reshare facts first,
which needs the rules to run again. Applying the rules makes new facts, which let the rules make more.

So the engine applies every rule, adds what it finds, and does it again,
until a round finds nothing new. The database is then at a \emph{fixpoint}
of the rules: applying any of them derives only facts it already has. That
is the whole of Datalog's meaning, in one sentence. The facts a program
implies are exactly those in the smallest database that contains its given
facts and is a fixpoint of its rules.

\section{Why it stops}

A loop that runs ``until nothing new appears'' needs an argument that it
stops. Here it is. Rules never invent constants: every constant in a
derived fact came from a given fact or from a rule. So there are only
finitely many facts that could ever be derived, a number of predicates
times the number of ways of filling their arguments with those constants.
And every round that does not stop adds at least one of them, while nothing
is ever removed. A finite supply that only grows must run out.

``Nothing is ever removed'' is something we can state. Every evaluator
promises it, together with the fixpoint itself, as the postcondition of
\code{run}. It is deferred in \code{Evaluator}, so every kind of evaluator
inherits both promises:

\begin{addbefore}{evaluate.nex}{to Evaluator, before solutions}{  -- Every set of bindings that satisfies the body of `r` in `db`.}
  -- How many times the rules were applied in turn.
  rounds: Integer

  run(rules: Array[Rule], db: Database)
  deferred
  ensure
    fixpoint: is_fixpoint(rules, db)
    only_grows: db.count >= old db.count
  end

  -- Note that another round of applying the rules has begun.
  count_round
  do
    rounds := rounds + 1
  end
\end{addbefore}

\code{count_round} is there because a subclass may read the fields it
inherits but not assign them, so it asks the class that declares
\code{rounds} to count.

\section{Checking the fixpoint}

The postcondition \code{fixpoint} is a strong check, and a simple one:
apply every rule once more, and see whether anything comes out that the
database lacks. It does not care how the evaluator got there, which is
exactly what we want from it. When a faster evaluator replaces the simple
one, this same check will say whether it is still right.

\begin{append}{evaluate.nex}
-- Does applying any of `rules` to `db` derive a fact it does not have?
function is_fixpoint(rules: Array[Rule], db: Database): Boolean
do
  result := true
  let e := create Naive_Evaluator.make
  across rules as r do
    across e.solutions(r, db) as b do
      if not db.contains(substitute(r.head, b)) then
        result := false
      end
    end
  end
end
\end{append}

\section{The naive evaluator}

The simplest evaluator does exactly what the definition says: apply every
rule to everything, round after round, until a round adds nothing. It is
called \emph{naive} evaluation, and the name is deserved, as
Chapter~\ref{ch:seminaive} will show.

\begin{append}{evaluate.nex}
-- Apply every rule to everything, again and again, until a round adds
-- nothing.
class Naive_Evaluator inherit Evaluator
create
  make
  do
  end
feature
  run(rules: Array[Rule], db: Database)
  do
    let changed := true
    from until not changed do
      changed := false
      count_round
      across rules as r do
        across solutions(r, db) as b do
          if db.add_derived(substitute(r.head, b), r, b) then
            changed := true
          end
        end
      end
    end
  end
end
\end{append}

The solutions of a rule are all found before any of its facts are added,
so a rule never reads facts it is adding in the middle of its own join.

\section{Trying it out}

\begin{tryit}{try_naive.nex}
intern Parser
intern io/Path
intern Evaluate

match parse_program(create Path.make(program.working_directory + "/social.dl").read_text) of
  Ok as ok then
    do
      let db := create Database.make
      across ok.value.facts as f do
        db.add(f)
      end
      let e := create Naive_Evaluator.make
      e.run(ok.value.rules, db)
      across db.facts_of("sees") as f do
        console.print_line(f.text)
      end
      console.print_line(db.count.to_string + " facts after " + e.rounds.to_string
        + " rounds, " + e.attempts.to_string + " facts looked at")
    end
  Err as e then console.print_line(e.error)
end
\end{tryit}
\tryoutput{try_naive}

The people each person follows come first, from the first rule in the
first round, then the longer chains, one reshare further each round. The last round
found nothing new, which is how the evaluator knew to stop, and the
postcondition then checked, independently, that nothing more follows.
